\documentclass[11pt]{article}
\usepackage{preamble}
\setlength{\parskip}{4pt}

\begin{document}

\daybanner{AP Statistics \textperiodcentered\ Unit 2 \textperiodcentered\ Topic 2.3 \textperiodcentered\ B: 2026-10-19 \textperiodcentered\ E: 2026-11-02 \textperiodcentered\ TEACHER KEY}{Can Chance Make a Streak?}

\sectionbanner{CED TETHER}

{\small
\begin{itemize}[leftmargin=1.5em,itemsep=2pt,topsep=2pt]
  \item \textbf{VAR-1.F} Identify questions suggested by patterns in data. \textbf{[Skill 1.A]}
  \item \textbf{VAR-1.F.1} Patterns in data do not necessarily mean that variation is not random.
\end{itemize}
}
\textbf{Essential question:} How can simulation distinguish a surprising-looking random pattern from evidence against a chance model?\par
\textbf{DOK-3 skill:} 4.A \textperiodcentered\ \textbf{FRQ pattern:} {\small\texttt{is-\allowbreak{}this-\allowbreak{}pattern-\allowbreak{}random-\allowbreak{}design-\allowbreak{}the-\allowbreak{}simulation}}\par
\textbf{DOK rationale:} Part (c) weighs competing claims using the day's numerical or graphical evidence and explains a limitation or consequence; the judgment requires linked reasoning beyond a calculation.\par

\frameworkphaseheader{Do Now (first take) $\rightarrow$ Explore (video + follow-along) $\rightarrow$ Exit (finish + turn in)}{1 $\rightarrow$ 3}{45}{%
\begin{itemize}[leftmargin=1.5em,itemsep=2pt,topsep=2pt]
  \item Show the board at the bell and collect a one-sentence first take without confirming a claim.
  \item Run the named video follow-along, then allow about 10 minutes for parts (a)--(c).
  \item Ask for evidence linked to a claim. Collect the sheet and hand-score part (c) with E/P/I.
\end{itemize}
}{%
\begin{itemize}[leftmargin=1.5em,itemsep=2pt,topsep=2pt]
  \item Commit to one claim and cite the sequence or dotplot before the lesson.
  \item Complete the follow-along about random patterns and simulation.
  \item Finish (a)--(c), revise the first take in the exit line, and turn in the sheet.
\end{itemize}
}{%
\begin{itemize}[leftmargin=1.5em,itemsep=2pt,topsep=2pt]
  \item Which number or visible feature supports your claim?
  \item Why does the other claim go beyond that evidence?
  \item What would you tell someone who chose the other claim?
\end{itemize}
}{Circulate and ask students to connect their choice to evidence. Score part (c) using the three required reasoning elements; do not require dropped extension work.}

\begin{callout}{\IconWarn\ WATCH FOR}{calloutred}
\begin{itemize}[leftmargin=1.5em,itemsep=2pt,topsep=2pt]
  \item A choice with copied numbers but no explanation of how they support it.
  \item Treating a model or average as a guaranteed individual outcome.
\end{itemize}
\end{callout}

\sectionbanner{SCORING PART (c) --- E / P / I}

\textbf{Expected elements}
\begin{itemize}[leftmargin=1.5em,itemsep=2pt,topsep=2pt]
  \item \textbf{reasoning-1} (required) --- Chooses B with 7 of 100 and avoids claiming proof of fairness
  \item \textbf{reasoning-2} (required) --- Matches one trial to 10 independent fair-coin flips and counts a longest run of at least 6
  \item \textbf{reasoning-3} (required) --- Explains why many repetitions improve the estimate
\end{itemize}

{\small\renewcommand{\arraystretch}{1.25}
\noindent\begin{tabular}{|p{0.5in}|p{5.9in}|}
\hline
\textbf{E} & Includes all three required reasoning elements above, with correct contextual evidence. \\ \hline
\textbf{P} & Includes at least one substantive correct reasoning link above, but omits or misstates another required element. \\ \hline
\textbf{I} & Includes no substantive correct reasoning link above; an unsupported choice or copied number alone is insufficient. \\ \hline
\end{tabular}}

\textbf{Common mistakes}
\begin{itemize}[leftmargin=1.5em,itemsep=2pt,topsep=2pt]
  \item Treating a run of six as impossible because each individual sequence is unlikely
  \item Using six flips rather than reproducing all 10 flips in one trial
  \item Counting only runs of exactly 6 and discarding longer runs
  \item Claiming that 7 percent proves the coin is fair
\end{itemize}

\clearpage
\focusbanner{TODAY'S PROBLEM --- annotated}

Suppose a coin advertised as fair is flipped 10 times and gives HHHHHHTTTT. To investigate what chance alone could produce, a class simulates 100 sets of 10 independent fair-coin flips. For each set, the class records the longest run of matching outcomes; the frequency table shows the results.\par\smallskip \textbf{Claim A:} ``The coin is rigged---six heads in a row cannot be chance.''\par \textbf{Claim B:} ``Streaks like this can happen by chance; simulation can tell us how unusual this one is.''\par
\begin{center}
\textbf{100 sets of 10 fair-coin flips}\par\smallskip
\begin{tabular}{|l|c|}
\hline
\textbf{Longest run} & \textbf{Sets} \\ \hline
\textbf{2} & 22 \\ \hline
\textbf{3} & 32 \\ \hline
\textbf{4} & 25 \\ \hline
\textbf{5} & 14 \\ \hline
\textbf{6} & 7 \\ \hline
\end{tabular}
\end{center}
\begin{firsttakebox}{FIRST TAKE (5 min)}
Which claim seems better supported right now? Cite one detail from the sequence or table.\par
\answer{Not graded. Preserve surprise; the key move is replacing visual intuition with a matched chance model.}
\end{firsttakebox}

\textit{Rules callout ``MODEL THE PATTERN, NOT THE STORY (keep on the page)'' is printed on the student sheet.}\par\medskip

\dokbadge{1}~\textbf{(a)} State the longest run in the observed sequence and the most common longest run in the simulated sets.\par
\answer{The observed sequence has a longest run of 6, the six consecutive heads. In the simulated sets, 3 is the most common longest-run value, appearing 32 times.}

\dokbadge{2}~\textbf{(b)} Estimate the chance of a run at least as long as the observed one. Use the number of simulated sets with such a run divided by 100.\par
\answer{Seven of the 100 simulated sets have a longest run of at least 6, so the estimate is $7/100=0.07$. In context, about 7 percent of sets of 10 independent flips from a fair coin produced a run at least as long as the observed run in this simulation.}

\dokbadge{3}~\textbf{(c)} Which claim is better supported? Use the simulation count and explain why it does not prove the coin is fair. To improve the estimate, describe one trial, what you would count, and why you would repeat many trials. Write 3--4 sentences.\par
\answer{Claim B is better supported: 7 of 100 fair-coin trials had a run at least 6 long, so the pattern can occur by chance. This result does not prove the actual coin is fair. One trial should use 10 independent fair-coin flips, recording the longest run of either outcome and counting a success when it is at least 6. Repeat many trials, such as 10,000, to reduce the chance variation in the estimated proportion.}

\begin{summaryexitbox}{EXIT}
One thing the video changed about my first take: \hrulefill
\end{summaryexitbox}

\end{document}
